\begingroup
\small\setstretch{1.12}\setlength{\tabcolsep}{3pt}
\setlength{\LTleft}{\fill}\setlength{\LTright}{\fill}
\begin{longtable}{@{}>{\raggedright\arraybackslash}p{0.15\linewidth}>{\raggedright\arraybackslash}p{0.15\linewidth}>{\raggedright\arraybackslash}p{0.21\linewidth}>{\raggedright\arraybackslash}p{0.21\linewidth}>{\raggedright\arraybackslash}p{0.21\linewidth}@{}}
\caption{實驗一不同靜態成員數之候選平均結果}\label{tab:results-static-size}\\
\toprule
模型數 & 候選數 & AUC & F1-score & Recall \\
\midrule\endfirsthead
\multicolumn{5}{c}{表~\thetable（續）}\\
\toprule
模型數 & 候選數 & AUC & F1-score & Recall \\
\midrule\endhead
\midrule\multicolumn{5}{r}{續下頁}\\\endfoot
\bottomrule\endlastfoot
2 & 9 & 0.7701 & 0.1763 & 0.2348 \\
3 & 18 & 0.7731 & 0.1745 & 0.2468 \\
4 & 15 & 0.7760 & 0.1752 & 0.2440 \\
5 & 6 & 0.7792 & 0.1763 & 0.2469 \\
6 & 1 & \textbf{0.7815} & \textbf{0.1815} & \textbf{0.2690} \\
\end{longtable}
\noindent{\footnotesize 資料來源：本研究實驗結果整理。\par 每個種子先平均同一成員數全部候選的測試分數，再取十個種子平均。粗體為各成員數中該指標最高值，僅供比較，不依此選成員數，也不表示顯著差異。本表沒有把候選再次組成新模型。}
\endgroup
